\documentclass[11pt]{article}
\usepackage[a4paper,margin=1in]{geometry}
\usepackage[T1]{fontenc}
\usepackage{lmodern}
\usepackage{microtype}
\usepackage{amsmath,amssymb,amsthm,mathtools}
\usepackage{booktabs}
\usepackage{xcolor}
\usepackage[numbers]{natbib}
\usepackage[colorlinks=true,allcolors=blue]{hyperref}

\newtheorem{theorem}{Theorem}
\newtheorem{proposition}[theorem]{Proposition}
\newtheorem{example}[theorem]{Example}
\newcommand{\Tr}{\operatorname{Tr}}
\newcommand{\Schatten}[1]{\mathcal S_{#1}}

\title{Orthogonal-Quotient Superloop Compression}
\author{Bin Wang}
\date{July 2026}

\begin{document}
\maketitle

\begin{abstract}
Cloud-extracted traces were represented in RH-242 as a graded difference of
physical loops and finite atomic counterloops.  We give an exact
cancellation-preserving grouping.  The selected generalized root space is
invariant, so its norm-one orthogonal quotient compression has power traces
equal to the cloud-extracted traces.  This avoids estimating the badly
conditioned Riesz projection.  A finite ordered-Schur audit verifies the
identity on 17 tractable archived endpoints and shows that one-step quotient
contraction still fails, motivating a block-power criterion.
\end{abstract}

\section{Invariant root space and orthogonal quotient}

Let $A\in\Schatten{2}(H)$ and let $\mathcal S$ be a finite algebraic multiset
of isolated nonzero eigenvalues of $A$.  Let $E$ be the sum of their full
generalized root spaces.  Then $E$ is finite dimensional and $A$-invariant.
Write $\Pi$ for the orthogonal projection onto $E$, $Q=I-\Pi$, and
\begin{equation}\label{eq:C}
 C=QAQ\big|_{E^\perp}:E^\perp\longrightarrow E^\perp.
\end{equation}
The complement $E^\perp$ need not be $A$-invariant: $C$ is the orthogonal
realization of the quotient operator, not a reducing restriction.

\begin{theorem}[orthogonal quotient trace identity]\label{thm:quotient}
For every integer $n\ge2$,
\begin{equation}\label{eq:identity}
 \Tr A^n-\sum_{s\in\mathcal S}s^n=\Tr C^n.
\end{equation}
The projection used in \eqref{eq:C} satisfies $\|\Pi\|=1$ whenever
$E\ne\{0\}$; no norm bound for the oblique Riesz projection is required.
\end{theorem}

\begin{proof}
Relative to the orthogonal decomposition $H=E\oplus E^\perp$, invariance of
$E$ gives
\[
 A=\begin{pmatrix}A_E&X\\0&C\end{pmatrix}.
\]
Every power has the same block upper-triangular form and diagonal blocks
$A_E^n$ and $C^n$.  Because $A\in\Schatten{2}$, these powers are trace class
for $n\ge2$.  Hence
$\Tr A^n=\operatorname{tr}A_E^n+\Tr C^n$.  The finite-dimensional trace of
$A_E^n$ is the algebraic power sum $\sum_{s\in\mathcal S}s^n$, including
multiplicity, which proves \eqref{eq:identity}.  The orthogonal projection
norm is elementary.
\end{proof}

This is precisely the RH-242 superloop cancellation \citep{WangRH242}
performed before taking absolute values.  In an ordered Schur basis, closed
block walks cannot contain the one-way coupling $X$: a walk that leaves the
quotient block cannot return.  The selected-block loops cancel the atomic
counterloops, leaving only quotient loops.

\begin{proposition}[quotient-kernel loops]\label{prop:kernel}
If $H=L^2(X,\mu)$ and $A$ has a Hilbert--Schmidt kernel, then $B=QAQ$ has a
Hilbert--Schmidt kernel $b$.  For every $n\ge2$,
\[
 \Tr C^n=\Tr B^n
 =\int_{X^n}b(x_0,x_1)b(x_1,x_2)\cdots
 b(x_{n-1},x_0)\,d\mu^n.
\]
Thus the grouped representation remains a genuine signed or complex
periodic-loop formula.
\end{proposition}

\begin{proof}
The Hilbert--Schmidt ideal is stable under bounded left and right
multiplication, so $B$ is Hilbert--Schmidt and vanishes on $E$ in the
orthogonal block form.  The standard trace formula for products of
Hilbert--Schmidt kernels gives the cyclic integral; its nonzero block is
$C^n$.
\end{proof}

\section{Why this differs from a Riesz estimate}

RH-232 found enormous Euclidean Riesz projection norms in the moving cloud
\citep{WangRH232}.  Orthogonal quotient compression removes that norm from
the fixed-operator trace identity, although it does not prove that the spaces
$E_\sigma$ vary uniformly.

\begin{example}[unbounded obliqueness, fixed quotient]\label{ex:model}
For distinct $\lambda,\mu$ let
\[
 A_M=\begin{pmatrix}\lambda&M\\0&\mu\end{pmatrix},
 \qquad E=\operatorname{span}\{e_1\}.
\]
The Riesz projection onto $E$ is
\[
 P_\lambda=\begin{pmatrix}1&M/(\lambda-\mu)\\0&0\end{pmatrix},
 \qquad
 \|P_\lambda\|=sqrt{1+|M/(\lambda-\mu)|^2},
\]
which diverges with $|M|$.  Yet $\Pi=\operatorname{diag}(1,0)$,
$C=[\mu]$, and
\[
 \operatorname{tr}A_M^n-\lambda^n=\mu^n
\]
for every $n$, independently of $M$.  Orthogonal grouping deletes purely
cross-block nonnormal mass, but says nothing about nonnormality internal to
$C$.
\end{example}

\section{Finite ordered-Schur audit}

We reconstruct every RH-222 endpoint of dimension at most 512
\citep{WangRH222}, Hardy-scale the matrix, and order its complex Schur form at
the archived radial gap.  The leading block contains Perron, parity, and the
selected cloud; the trailing block is $C$.  We compare its power traces with
the projection-free RH-236 archive \citep{WangRH236}.

\begin{center}
\small
\begin{tabular}{lr}
\toprule
quantity & value\\
\midrule
audited endpoints & 17\\
selected block dimensions & 6--25\\
rank mismatches & 0\\
maximum Schur trace-partition error & $3.99\times10^{-15}$\\
maximum RH-236 residual difference & $6.31\times10^{-12}$\\
quotient Frobenius square & 5.07--60.09\\
old-budget improvement factor & 1.20--1.62\\
quotient one-step norm & 1.277--2.346\\
first contractive power depth & 3--7\\
maximum $\|C^{12}\|$ & $1.370\times10^{-5}$\\
\bottomrule
\end{tabular}
\end{center}

The exact theorem explains why cross-block mass is irrelevant, but the
finite numbers also prevent an overclaim: every audited one-step quotient
norm exceeds one, and the Frobenius budget improves only modestly.  Deep
power contraction is striking finite evidence, not a uniform theorem.

\section{Boundary and next condition}

Theorem~\ref{thm:quotient} is exact at each fixed noise once the algebraic
selected subspace exists.  It does not provide a noise-uniform method to
identify that subspace, a uniform kernel bound, or a uniform contractive
power.  Only 17 of 32 archived endpoints were densely decomposed, and no
interval certification was attempted.  The next obligation is therefore an
explicit $m$-block trace-envelope theorem and a uniform certificate for its
hypotheses.  The anchored selector obstruction of RH-244 remains separate.

Gate A remains open and Gates B--E are untouched.  No Hilbert--P\'olya
operator, zeta-divisor equality, Riemann-zero identification, or RH
implication is claimed.

\bibliographystyle{plainnat}
\bibliography{references}
\end{document}
